\section{Evaluación final}

\subsection{Apertura única del holdout}

Una vez fijados la configuración de la cabeza, los modelos participantes y la
regla de voting, cada cabeza se reajustó con las 28\,418 imágenes de desarrollo.
El proceso comprobó que no existiera previamente
\texttt{outputs/etapa\_2/final/metrics.json}; después creó el directorio final,
evaluó las 5\,015 filas congeladas y registró
\texttt{evaluation\_count=1}. El JSON conserva además el hash del CSV de
holdout. Esta protección no es criptográfica frente a una persona que borre
archivos, pero deja auditable la ejecución usada por el informe.

\subsection{Resultados globales y por clase}

\begin{table}[H]
  \centering
  \caption{Métricas obtenidas en la única evaluación del holdout.}
  \label{tab:final-global}
  \resizebox{0.92\textwidth}{!}{\begin{tabular}{rrrrrr}
\toprule
\textbf{Accuracy} & \textbf{P-macro} & \textbf{R-macro} & \textbf{F1-macro} & \textbf{F1 pond.} & \textbf{N} \\
\midrule
0.9537 & 0.9012 & 0.9092 & 0.9035 & 0.9537 & 5015 \\
\bottomrule
\end{tabular}
}
\end{table}

\begin{table}[H]
  \centering
  \caption{Precision, recall y F1 finales para las nueve clases.}
  \label{tab:final-class}
  \small
  \begin{tabular}{lrrrr}
\toprule
\textbf{Clase} & \textbf{Soporte} & \textbf{Precision} & \textbf{Recall} & \textbf{F1} \\
\midrule
Roya común & 338 & 0.9970 & 0.9675 & 0.9820 \\
Gusano cogollero & 729 & 0.9517 & 0.9739 & 0.9627 \\
Mancha gris & 290 & 0.9027 & 0.9276 & 0.9150 \\
Sana & 1311 & 0.9674 & 0.9748 & 0.9711 \\
Necrosis letal & 962 & 0.9927 & 0.9896 & 0.9912 \\
Def. nitrógeno & 127 & 0.7386 & 0.8898 & 0.8071 \\
Tizón norteño & 1024 & 0.9653 & 0.9248 & 0.9446 \\
Def. fósforo & 141 & 0.8553 & 0.9220 & 0.8874 \\
Def. potasio & 93 & 0.7403 & 0.6129 & 0.6706 \\
\bottomrule
\end{tabular}

\end{table}

\begin{figure}[H]
  \centering
  \includegraphics[width=0.90\textwidth]{final_per_class.pdf}
  \caption{Rendimiento final por clase. El soporte desigual no cambia el peso
  de cada barra en macro-F1.}
  \label{fig:final-class}
\end{figure}

El ensemble obtuvo 0.9537 de accuracy y F1 ponderado, 0.9012 de precision
macro, 0.9092 de recall macro y 0.9035 de macro-F1. El resultado quedó 0.0062
por debajo de la media de validación cruzada y dentro de una desviación
estándar de ella. La validación fija había sido más optimista (0.9202), algo
esperable porque allí también se eligieron los pesos.

La diferencia entre accuracy y macro-F1 vuelve a concentrarse en tres clases.
Necrosis letal obtuvo F1 0.9912 y roya común 0.9820; sana, gusano cogollero y
tizón norteño también superaron 0.94. Potasio quedó en 0.6706, nitrógeno en
0.8071 y fósforo en 0.8874. El ensemble mantiene el rendimiento global, pero no
convierte las deficiencias en un problema resuelto.

\subsection{Matriz de confusión}

\begin{figure}[H]
  \centering
  \includegraphics[width=\textwidth]{final_confusion_matrices.pdf}
  \caption{Matriz final en conteos y normalizada por clase real. La segunda
  vista permite leer clases pequeñas sin que desaparezcan detrás de la clase
  sana.}
  \label{fig:final-confusion}
\end{figure}

Hubo 232 errores en 5,015 imágenes. Los dos flujos individuales más grandes
fueron potasio hacia nitrógeno (28 casos) y tizón norteño hacia sana (25).
También aparecieron 21 casos de tizón como mancha gris y 13 en el sentido
contrario. La matriz normalizada hace visible que los 36 fallos de potasio
pesan mucho más para su clase que los 77 fallos totales de tizón para una clase
con 1,024 ejemplos.

La diagonal de roya parece casi perfecta en conteos, pero esconde una
diferencia de entorno que se revisa en fairness: 327 de 338 casos fueron
correctos y los once errores se concentran proporcionalmente en las pocas
imágenes de campo. El problema no está en reconocer el patrón de laboratorio,
sino en sostenerlo cuando cambia la captura.

\subsection{Etapa 1 frente a Etapa 2}

\begin{table}[H]
  \centering
  \caption{Referencias de ambas etapas. Accuracy y macro-F1 no son una
  comparación controlada porque cambiaron corpus, división y entrenamiento.}
  \label{tab:e1-e2}
  \resizebox{\textwidth}{!}{\begin{tabular}{lrrl}
\toprule
\textbf{Métrica} & \textbf{Etapa 1} & \textbf{Etapa 2} & \textbf{Lectura} \\
\midrule
Accuracy & 0.9521 & 0.9537 & Corpus y protocolo distintos \\
Macro-F1 & 0.9146 & 0.9035 & Corpus y protocolo distintos \\
Recall potasio & -- & 0.6129 & No comparable \\
F1 potasio & 0.62 & 0.6706 & Referencia B0 histórica \\
Imágenes potasio & 266 & 621 & +355 \\
\bottomrule
\end{tabular}
}
\end{table}

La tabla no atribuye el cambio a Optuna ni a la ampliación. El contraste
controlado de tuning está en la Tabla~\ref{tab:baseline-tuned}; aquí la Etapa 1
aporta el antecedente del proyecto y la Etapa 2 aporta un holdout que sí
permaneció cerrado. Esta diferencia metodológica vale más que aparentar una
ganancia directa entre dos números obtenidos bajo condiciones distintas.

Como referencia, accuracy cambió de 0.9521 a 0.9537 y macro-F1 de 0.9146 a
0.9035. No se interpreta la caída de 0.0111 como empeoramiento causal: el
holdout nuevo contiene otro corpus y la cabeza no fue ajustada end-to-end. En
potasio, el F1 pasó de 0.62 a 0.6706 y el soporte evaluado de 40 a 93. La señal
es favorable, pero todavía insuficiente y tampoco aislable de los cambios de
datos y protocolo.
